%! Author = Omar Iskandarani
%! Date = 4/14/2026
%! Affiliation = Independent Researcher, Groningen, The Netherlands
%! ORCID = 0009-0006-1686-3961
%! DOI = 10.5281/zenodo.19589794

\newcommand{\paperdoi}{10.5281/zenodo.19589794}
\newcommand{\papertitle}{Delay-Coupled Phase Locking on Shared Circulation Carriers in a Minimal Toy Model\\
\large A Reduced Dynamical-Systems Closure for Delay, Coupling, and Coherence}
\newcommand{\paperkeywords}{entanglement-like correlations, toy model, delay-coupled phase dynamics, circulation overlap, coherence length, hydrodynamic-topological model}

\documentclass[11pt]{article}

%==============================================================================
% PACKAGES
%==============================================================================
\usepackage{lmodern}
\usepackage{microtype}
\usepackage{amsmath,amssymb,amsfonts,amsthm,bm}
\usepackage{physics}
\usepackage{pgfplots}
\pgfplotsset{compat=1.18}
\usepackage{tikz}
\usepackage{mathtools}
\usepackage{enumitem}
\usepackage{hyperref}
\usepackage[a4paper,margin=1in]{geometry}
\usepackage[absolute,overlay]{textpos}
\usepackage[T1]{fontenc}
\usepackage[utf8]{inputenc}

\usepackage[nameinlink,noabbrev]{cleveref}

%==============================================================================
% SST MACRO PRELUDE
%==============================================================================
\newcommand{\vswirl}{\mathbf{v}_{\!\boldsymbol{\circlearrowleft}}}
\newcommand{\vnorm}{\lVert \mathbf{v}_{\!\boldsymbol{\circlearrowleft}}\rVert}
\newcommand{\rhoF}{\rho_{\!f}}
\newcommand{\rhoE}{\rho_{\!E}}
\newcommand{\rhoM}{\rho_{\!m}}
\newcommand{\rc}{r_c}
\newcommand{\rhocore}{\rho_{\text{core}}}
\newcommand{\clockfield}{\chi}
\newcommand{\Gammazero}{\Gamma_0}
\newcommand{\SwirlClock}{S_t^{\boldsymbol{\circlearrowleft}}}


\begin{document}
    \thispagestyle{empty}%
    \begin{center}
        \Large \bfseries \papertitle \par \vspace{1cm}
        {\Large \itshape \textbf{Omar Iskandarani}\textsuperscript{\textbf{*}} \par} \vspace{0.5cm}
        {\small Independent Researcher, Groningen, The Netherlands\par} \vspace{0.5cm}
        {\today \par}  \vspace{0.5cm}
    \end{center}
    
    
    \begin{abstract}
        We formulate a minimal Swirl--String Theory (SST) toy model for delay-coupled phase locking on shared circulation carriers. The central proposal is that correlated subsystems may be modeled not as independent pointlike entities linked only by an abstract nonlocal state, but as localized readout sites on a common circulation carrier embedded in the swirl medium. At the reduced effective level, the dynamics are described by a delayed phase system with three macroscopic parameters: a transport delay \(\tau\), a phase-locking rate \(\kappa\), and a coherence length \(\ell_\tau\). We show that the canonical SST identification
        \[
            \omega_c=\frac{\vnorm}{\rc}
        \]
        immediately fixes the transport time scale and strongly constrains the remaining two parameters to the form
        \[
            \tau=\frac{L}{\vnorm},\qquad
            \kappa_{AB}=\omega_c\,\mathcal{O}_{AB},\qquad
            \ell_\tau=\rc\,\mathcal{Q},
        \]
        where \(\mathcal{O}_{AB}\) is a dimensionless shared-circulation overlap factor and \(\mathcal{Q}\) is a dimensionless phase-quality factor. The present article is intended as a dynamical-systems and reduced-modeling contribution rather than a quantum-foundations derivation: no claim is made here to a derivation of Bell/CHSH phenomenology, the Born rule, or a direct experimental fit to quantum-correlation data. We distinguish derived results, closure assumptions, and open problems, propose minimal falsifiers for each sector, and provide a reduced observable construction together with a worked coherence estimate based on topological leakage.
    \end{abstract}
    
    Keywords: \paperkeywords
    
    
    \begin{textblock*}{1\textwidth}(0.15\textwidth,1.1\textheight)\raggedright\footnotesize\renewcommand{\arraystretch}{1.0} \noindent\rule{\textwidth}{0.4pt}\\[0.5em]
        \textsuperscript{\textbf{*}} Email: \texttt{info@omariskandarani.com}\\
        ORCID: \texttt{\href{https://orcid.org/0009-0006-1686-3961}{0009-0006-1686-3961}}\\
        DOI: \href{https://doi.org/\paperdoi}{\paperdoi}
    \end{textblock*}
    
    \newpage
    
    
    \section{Introduction}
        
        The standard quantum formalism represents entanglement through nonfactorizable states in Hilbert space. While empirically successful, that description leaves open the ontological status of the correlation carrier. Hydrodynamic and stochastic reformulations of quantum structure have a substantial history, ranging from Madelung's continuum rewriting of the Schrödinger equation to Nelson's stochastic mechanics, generalized delayed phase-locking models \cite{Yeung1999}, Bohm's hidden-variable program, and modern analogue-gravity or superfluid-vacuum programs \cite{Madelung1927,Nelson1966,Bohm1952,Barcelo2011,Volovik2003}. The present article does not attempt to review that literature comprehensively. Instead, it isolates one specific reduced question within SST.

        In SST, the primitive ontology is different: the vacuum is modeled as an incompressible, inviscid structured medium, and localized matter-like excitations are topological swirl strings rather than point particles. Within such a framework, it is natural to ask whether nontrivial long-range correlations can be re-expressed, at least in a reduced model, as dynamical consequences of shared circulation structure, finite propagation, and phase-locking.
        
        The present article does not claim a completed derivation of Bell correlations from first principles. Rather, it isolates the minimum dynamical closure problem in the form of a toy model. If correlated subsystems \(A\) and \(B\) are treated as local phase-bearing sites on a common circulation carrier, then the reduced effective theory is determined by three scales:
        \begin{enumerate}[label=(\roman*)]
    \item a transport delay \(\tau\),
    \item a phase-locking rate \(\kappa\),
    \item a coherence length \(\ell_\tau\).
        \end{enumerate}
        Instead, it isolates a minimal mechanistic sector whose role is to determine whether SST admits a nontrivial reduced model for delay, geometric coupling, and finite coherence on shared circulation carriers.

        The paper is therefore positioned primarily as a dynamical-systems and reduced-modeling study. Throughout this article, the phrase ``entanglement-like'' is used only in a restricted operational sense: separated readout sites on a shared carrier may display angle-dependent correlations not attributable to independent local phases, without thereby claiming a completed reproduction of quantum-mechanical entanglement. In particular, the present manuscript does not derive Bell/CHSH inequalities, the Born rule, or a direct comparison with experimental quantum-correlation datasets.

        The main result of this paper is that SST already fixes the parametric structure of all three at the level of this toy model:
        \[
            \tau=\frac{L}{\vnorm},\qquad
            \kappa_{AB}=\omega_c\,\mathcal{O}_{AB},\qquad
            \ell_\tau=\rc\,\mathcal{Q}.
        \]
        What is specific to SST in this reduction is not merely the presence of a characteristic frequency, but the fact that the transport, coupling, and coherence sectors are all tied back to the same primitive core scales \((\rc,\vnorm)\) through the identifications \(\omega_c=\vnorm/\rc\), \(\kappa_{AB}=\omega_c\mathcal{O}_{AB}\), and \(\ell_\tau=\rc\mathcal{Q}\).

        The only remaining physical inputs are geometric overlap and the mechanism of phase-memory loss.

        The purpose of the paper is therefore not to replace the quantum formalism, but to ask whether a minimal hydrodynamic-topological model can isolate a plausible SST mechanism for transport delay, geometric coupling, and finite coherence on shared circulation carriers. The main output is thus a reduced closure model, not yet a quantum-foundations account.
    
    \section{Epistemic status of the claims}
        
        To keep the logical status explicit, we separate the statements as follows.
        
        \begin{itemize}
            \item \textbf{Derived from current SST anchors:}
            \[
                \omega_c=\frac{\vnorm}{\rc},\qquad
                \tau=\frac{L}{\vnorm}.
            \]
            \item \textbf{Constrained toy-model closure forms:}
            \[
                \kappa_{AB}=\omega_c\,\mathcal{O}_{AB},\qquad
                \ell_\tau=\rc\,\mathcal{Q}.
            \]
            \item \textbf{Modeled but not yet first-principles derived:}
            explicit formulas for \(\mathcal{O}_{AB}\) and \(\mathcal{Q}\).
            \item \textbf{Open program:}
            recovery of exact Bell/CHSH statistics and operational no-signaling from the underlying SST medium dynamics.
            \item \textbf{Not claimed here:}
            equivalence to standard quantum mechanics, a derivation of the Born rule, a Bell/CHSH derivation, a direct experimental comparison to measured quantum-correlation data, or a full SST account of quantum measurement.
        \end{itemize}
    
    \section{Canonical SST anchors}
        
        \paragraph{Notation convention.}
        Throughout this article, the characteristic swirl speed is denoted exclusively by the vector symbol \(\mathbf{v}_{\!\boldsymbol{\circlearrowleft}}\), with magnitude \(\lVert \mathbf{v}_{\!\boldsymbol{\circlearrowleft}}\rVert\). No alternative symbols are intended for this quantity. This convention is enforced in all formulas below to avoid ambiguity.

        For the purposes of the present toy model, SST is taken to supply three working primitives: a core radius \(\rc\), a characteristic swirl speed \(\vnorm\), and a fundamental circulation scale. These are adopted here as model inputs rather than derived within the present article. In the broader SST program, these quantities are calibrated to the electron/Compton/core closure chain; in the present paper they function only as reduced-model constants.

        Their intended provenance may be summarized briefly as follows. SST adopts the classical electron radius
        \begin{equation}
            r_e=\frac{\alpha\hbar}{m_e c},
        \end{equation}
        and takes the core radius to be
        \begin{equation}
            \rc=\frac{r_e}{2}.
        \end{equation}
        The characteristic angular rate is identified with the electron Compton angular frequency
        \begin{equation}
            \omega_c=\frac{m_e c^2}{\hbar},
        \end{equation}
        and the characteristic swirl speed is then fixed by the kinematic closure
        \begin{equation}
            \vnorm=\omega_c \rc.
        \end{equation}
        The present article does not defend that broader closure program; it uses only the resulting reduced constants.

        The fundamental circulation is
        \begin{equation}
            \Gammazero=2\pi \rc \vnorm.
        \end{equation}
        The canonical Compton-frequency identification is
        \begin{equation}
            \omega_c=\frac{\vnorm}{\rc}.
            \label{eq:omega_c}
        \end{equation}
        which is dimensionally exact:
        \[
            [\omega_c]=\frac{\mathrm{m\,s^{-1}}}{\mathrm{m}}=\mathrm{s^{-1}}.
        \]
        
        Using the canonical constants
        \[
            \vnorm=1.09384563\times 10^6\ \mathrm{m\,s^{-1}},
            \qquad
            \rc=1.40897017\times 10^{-15}\ \mathrm{m},
        \]
        one obtains
        \begin{equation}
            \omega_c
            =
            \frac{1.09384563\times 10^6}{1.40897017\times 10^{-15}}
            \approx 7.76344\times 10^{20}\ \mathrm{s^{-1}}.
            \label{eq:omega_c_num}
        \end{equation}
        
        This identification is crucial because it fixes the intrinsic turnover time of an elementary circulation core within the toy model:
        \begin{equation}
            \omega_c^{-1}\approx 1.288\times 10^{-21}\ \mathrm{s}.
        \end{equation}
    
    \section{Effective delayed phase model}
        
        We model two correlated sites \(A\) and \(B\) by phase variables \(\phi_A(t)\) and \(\phi_B(t)\) satisfying the symmetric delayed phase system, written in the standard spirit of delayed Kuramoto-type phase locking \cite{Yeung1999},
        \begin{equation}
            \dot{\phi}_A(t)=\omega_0+\kappa_{AB} \sin\!\big[\phi_B(t-\tau)-\phi_A(t)\big],
            \label{eq:phiA}
        \end{equation}
        \begin{equation}
            \dot{\phi}_B(t)=\omega_0+\kappa_{AB} \sin\!\big[\phi_A(t-\tau)-\phi_B(t)\big].
            \label{eq:phiB}
        \end{equation}
        The structure of \cref{eq:phiA,eq:phiB} is intentionally orthodox: it is a finite-delay phase-locking model. The SST content lies not in the existence of such equations, but in the interpretation of \(\tau\), \(\kappa\), and \(\ell_\tau\) as hydrodynamic-topological quantities.
        
        Define the phase difference
        \begin{equation}
            \Delta(t)=\phi_A(t)-\phi_B(t).
        \end{equation}
        Subtracting \cref{eq:phiA,eq:phiB} yields
        \begin{equation}
            \dot{\Delta}(t)
            =
            -2\kappa_{AB} \sin\!\big(\Delta(t-\tau)\big).
            \label{eq:deltaeq}
        \end{equation}
        Stationary branches satisfy
        \begin{equation}
            \sin(\Delta^\ast)=0
            \quad\Rightarrow\quad
            \Delta^\ast=n\pi,\qquad n\in\mathbb{Z}.
            \label{eq:stationary_branch}
        \end{equation}
        Thus a delay-coupled closed phase system admits discrete branch selection even before any quantum-statistical interpretation is imposed.

    \section{Reduced observables and correlation function}

        To make the toy model operational, we define binary readout observables by thresholding the local phases against analyzer angles \(\theta_A\) and \(\theta_B\):
        \begin{equation}
            A(\theta_A,t)=\operatorname{sign}\!\big[\cos(\phi_A(t)-\theta_A)\big],
            \label{eq:Aobs}
        \end{equation}
        \begin{equation}
            B(\theta_B,t)=\operatorname{sign}\!\big[\cos(\phi_B(t)-\theta_B)\big].
            \label{eq:Bobs}
        \end{equation}
        These observables take values in \(\{-1,+1\}\) and define the reduced-model correlation function
        \begin{equation}
            C(\theta_A,\theta_B)
            :=
            \left\langle
            A(\theta_A,t)\,B(\theta_B,t)
            \right\rangle_t,
            \label{eq:Cobs}
        \end{equation}
        where \(\langle\cdot\rangle_t\) denotes a long-time average on a stationary branch or an ensemble average over initial conditions within the toy model.

        Equations \cref{eq:Aobs,eq:Bobs,eq:Cobs} do not constitute a Bell/CHSH framework by themselves. Their purpose is narrower: they provide an explicit observable layer showing how the delayed phase system can generate nontrivial angle-dependent correlations in a self-contained reduced setting.
    
    \section{Derivation of the transport delay}
        
        Consider a phase disturbance propagating along a circulation structure of arclength \(L\). If the transport speed is the characteristic SST swirl speed \(\vnorm\), then the delay is simply
        \begin{equation}
            \boxed{
                \tau=\frac{L}{\vnorm}.
            }
            \label{eq:tau_basic}
        \end{equation}
        This is the natural kinematic first-pass closure consistent with the canonical identification \(\omega_c=\vnorm/\rc\).
        
        Dimensional check:
        \[
            [\tau]=\frac{\mathrm{m}}{\mathrm{m\,s^{-1}}}=\mathrm{s}.
        \]
        
        Equivalently,
        \begin{equation}
            \tau\,\omega_c
            =
            \frac{L}{\vnorm}\frac{\vnorm}{\rc}
            =
            \frac{L}{\rc}.
            \label{eq:tau_omega}
        \end{equation}
        Thus the delay measured in carrier cycles is simply the geometric loop length in units of the core radius.
        
        \subsection*{Interpretation}
            Equation \eqref{eq:tau_omega} is important. It says that a circulation loop of length \(L\) contains \(L/\rc\) core-turnover lengths, so the delay is not an independent free time parameter but a geometric bookkeeping of how many core scales fit around the carrier.
    
    \section{Why the coupling rate must be geometric}
        
        In \cref{eq:phiA,eq:phiB}, \(\kappa\) multiplies a sine and must therefore have units of inverse time:
        \[
            [\kappa]=\mathrm{s^{-1}}.
        \]
        Since \(\omega_c\) is the only intrinsic SST primitive with those units, \(\kappa\) must take the form
        \begin{equation}
            \kappa=\omega_c \times (\text{dimensionless factor}).
            \label{eq:kappa_structure}
        \end{equation}
        That dimensionless factor cannot be fixed by \(\rhoF\) alone. Density sets energetic scale, but not the degree to which two localized phase probes sample the same circulation object. Therefore \(\kappa\) must depend on geometry.

        In the present SST reading, this geometric dependence should not be interpreted as mere spatial proximity. A useful physical distinction is between vector-like and scalar-like sectors of the carrier. Local circulation direction and local phase gradients are vector-like quantities and may partially cancel under mirrored or anti-oriented configurations, whereas overlap strength and stored phase-locking capacity behave as scalar-like quantities that may add constructively. The reduced coupling factor \(\mathcal{O}_{AB}\) is intended to capture this scalar compatibility of the shared carrier rather than only naive co-location.
        
        We introduce the \emph{shared-circulation overlap factor}
        \begin{equation}
            \mathcal{O}_{AB}\in[-1,1],
        \end{equation}
        and write the general SST closure
        \begin{equation}
            \boxed{
                \kappa_{AB}=\omega_c\,\mathcal{O}_{AB}.
            }
            \label{eq:kappa_general}
        \end{equation}
    
    \section{Geometric derivation of the overlap factor}
        
        The heuristic closure
        \[
            \mathcal{O}_{AB}\approx \frac{\ell_{\mathrm{shared}}}{L}
        \]
        can be strengthened by deriving the overlap factor directly from the SST interaction energy of two circulation carriers.
        
        Let \(\mathbf{u}_A(\mathbf{x})\) and \(\mathbf{u}_B(\mathbf{x})\) denote the velocity fields induced by two localized phase-carrying swirl structures \(A\) and \(B\). The natural SST interaction energy is
        \begin{equation}
            E_{AB}^{\mathrm{int}}
            =
            \rho_{\!f}\int_{\mathbb{R}^3}
            \mathbf{u}_A(\mathbf{x})\cdot \mathbf{u}_B(\mathbf{x})\,d^3x,
            \label{eq:EABint}
        \end{equation}
        with self-energies
        \begin{equation}
            E_A
            =
            \rho_{\!f}\int_{\mathbb{R}^3}\|\mathbf{u}_A(\mathbf{x})\|^2\,d^3x,
            \qquad
            E_B
            =
            \rho_{\!f}\int_{\mathbb{R}^3}\|\mathbf{u}_B(\mathbf{x})\|^2\,d^3x.
            \label{eq:EAEB}
        \end{equation}
        
        This suggests the normalized overlap functional
        \begin{equation}
            \boxed{
                \mathcal{O}_{AB}
                :=
                \frac{\displaystyle \int_{\mathbb{R}^3}\mathbf{u}_A\cdot \mathbf{u}_B\,d^3x}
                {\displaystyle
                \sqrt{
                    \left(\int_{\mathbb{R}^3}\|\mathbf{u}_A\|^2\,d^3x\right)
                    \left(\int_{\mathbb{R}^3}\|\mathbf{u}_B\|^2\,d^3x\right)
                }}
            }
            \label{eq:OAB_energy}
        \end{equation}
        which is dimensionless and bounded by
        \begin{equation}
            -1 \le \mathcal{O}_{AB}\le 1
        \end{equation}
        by the Cauchy--Schwarz inequality.

        \paragraph{Transparency of the reduction chain.}
        The approximation hierarchy used in this section is as follows. At the canonical level, \(\mathcal{O}_{AB}\) is defined by the normalized cross-energy overlap \cref{eq:OAB_energy}. The reduction to \cref{eq:OAB_shared_general} assumes thin carriers with identical transverse profiles and aligned tangent directions on the shared segment. The final form \(\ell_{\mathrm{shared}}/L\) is the equal-length specialization \(L_A=L_B\). If these assumptions fail, the canonical overlap definition remains valid, but the simplified shared-length reduction no longer need hold.

        It is important that this hierarchy be read as a progression from a field-based definition to increasingly specialized geometric limits. The simplest shared-length picture is therefore not the primary SST object, but the most compressed reduction of a more general carrier-compatibility functional. In particular, once tangent alignment, chirality, or local phase transport differ appreciably between two sectors, the overlap must be understood as a genuine field overlap rather than a pure arclength fraction.
        
        The coupling rate is then
        \begin{equation}
            \boxed{
                \kappa_{AB}=\omega_c\,\mathcal{O}_{AB}.
            }
            \label{eq:kappa_energy_overlap}
        \end{equation}
        
        \subsection{Thin-filament reduction}
            
            To connect \cref{eq:OAB_energy} with the simpler shared-length closure, assume both carriers are thin filaments with the same transverse core profile \(f(\mathbf{x}_\perp)\), equal characteristic amplitude \(u_0\), and aligned tangent direction \(\hat{\mathbf{t}}(s)\) on the shared segment:
            \begin{equation}
                \mathbf{u}_A(\mathbf{x})
                \approx
                u_0\,f(\mathbf{x}_\perp)\,\chi_A(s)\,\hat{\mathbf{t}}(s),
                \qquad
                \mathbf{u}_B(\mathbf{x})
                \approx
                u_0\,f(\mathbf{x}_\perp)\,\chi_B(s)\,\hat{\mathbf{t}}(s),
            \end{equation}
            where \(\chi_A,\chi_B\) are support functions along arc length \(s\).
            
            Then
            \begin{equation}
                \int \mathbf{u}_A\cdot \mathbf{u}_B\,d^3x
                =
                u_0^2
                \left(\int |f|^2\,d^2x_\perp\right)
                \left(\int \chi_A(s)\chi_B(s)\,ds\right),
            \end{equation}
            while
            \begin{equation}
                \int \|\mathbf{u}_A\|^2\,d^3x
                =
                u_0^2
                \left(\int |f|^2\,d^2x_\perp\right)L_A,
            \end{equation}
            \begin{equation}
                \int \|\mathbf{u}_B\|^2\,d^3x
                =
                u_0^2
                \left(\int |f|^2\,d^2x_\perp\right)L_B.
            \end{equation}
            
            Therefore
            \begin{equation}
                \boxed{
                    \mathcal{O}_{AB}
                    =
                    \frac{\ell_{\mathrm{shared}}}{\sqrt{L_A L_B}},
                }
                \label{eq:OAB_shared_general}
            \end{equation}
            where
            \[
                \ell_{\mathrm{shared}}=\int \chi_A(s)\chi_B(s)\,ds.
            \]
            
            For equal loop lengths \(L_A=L_B=L\), \cref{eq:OAB_shared_general} reduces to
            \begin{equation}
                \boxed{
                    \mathcal{O}_{AB}
                    =
                    \frac{\ell_{\mathrm{shared}}}{L}.
                }
                \label{eq:OAB_shared_equal}
            \end{equation}
            
            Thus the earlier shared-length law is not a separate ansatz but the equal-length thin-filament limit of the energy-normalized SST overlap functional.
        
        \subsection{Sign and chirality}
            
            Because \(\mathcal{O}_{AB}\) is built from \(\mathbf{u}_A\cdot \mathbf{u}_B\), it is sensitive to orientation:
            
            \[
                \mathcal{O}_{AB}>0
                \quad\text{for co-oriented circulation,}
            \]
            \[
                \mathcal{O}_{AB}<0
                \quad\text{for anti-oriented circulation,}
            \]
            \[
                \mathcal{O}_{AB}\approx 0
                \quad\text{for orthogonal or disjoint carriers.}
            \]
            
            Accordingly, the effective phase equation may be written in the signed form
            \begin{equation}
                \dot{\phi}_A
                =
                \omega_0
                +
                \omega_c\,\mathcal{O}_{AB}
                \sin[\phi_B(t-\tau)-\phi_A(t)],
            \end{equation}
            with an analogous equation for \(\phi_B\). This retains the chirality-sensitive structure of the overlap rather than absorbing it into a positive scalar coupling.
        
        \subsection{Physical reading}
            Equation \eqref{eq:kappa_energy_overlap} implies that correlation strength is not primitive. In the toy model, it is controlled by the normalized overlap of genuinely shared circulation structure. In SST language, two subsystems are strongly correlated not because they are assigned a mysterious abstract relation, but because they are modeled as two local manifestations of the same phase-carrying object.

            This also suggests a refinement of interpretation: \(\mathcal{O}_{AB}\) should be read not only as a measure of shared support in space, but more generally as a reduced measure of geometric and kinematic compatibility on a common carrier. At minimum this includes shared support, tangent alignment, and orientation; in more refined versions it may also include resonant compatibility of the local phase mode. The present article does not yet build such resonance weighting explicitly, but the overlap formalism is intended to leave room for that stronger interpretation.

        \subsection{A minimal explicit example}
            
            The overlap functional \(\mathcal{O}_{AB}\) becomes concrete once explicit carrier fields are chosen. As the simplest reduced example, consider two identical coaxial thin carriers with the same transverse profile \(f(\mathbf{x}_\perp)\), the same characteristic amplitude \(u_0\), equal arclength \(L\), and fully overlapping support along the carrier coordinate \(s\). Then \(\chi_A(s)=\chi_B(s)\), so \(\ell_{\mathrm{shared}}=L\), and \cref{eq:OAB_shared_equal} gives
            \begin{equation}
                \mathcal{O}_{AB}=1
            \end{equation}
            for co-oriented circulation.

            If the same two carriers are anti-oriented, the sign of \(\mathbf{u}_A\cdot\mathbf{u}_B\) reverses everywhere in the overlap region and one obtains
            \begin{equation}
                \mathcal{O}_{AB}=-1.
            \end{equation}
            In the opposite limit of disjoint or ideally orthogonal carriers, the overlap becomes small and
            \begin{equation}
                \mathcal{O}_{AB}\approx 0.
            \end{equation}

            A slightly less degenerate intermediate case is two equal-length coaxial carriers with only partial shared support. If the common support occupies a fraction \(f\in(0,1)\) of the total arclength, so that
            \[
                \ell_{\mathrm{shared}}=fL,
            \]
            then \cref{eq:OAB_shared_equal} gives
            \begin{equation}
                \mathcal{O}_{AB}=f.
            \end{equation}

            These examples are intentionally minimal. They do not yet constitute a full Biot--Savart evaluation for realistic SST knot geometries. Their role is narrower: they show how the abstract overlap functional maps onto the simplest carrier configurations and how its sign tracks co- versus anti-oriented circulation.

        \subsection{Expected bounds beyond the parallel limit}
            
            The usefulness of \(\mathcal{O}_{AB}\) is not restricted to the degenerate parallel-filament limit. Even before carrying out a full Biot--Savart evaluation for realistic SST knots, one can state simple symmetry-based bounds for nontrivial configurations. For a symmetric Hopf-link-style pair of linked but approximately orthogonal carriers, the geometric linkage is nonzero while the local velocity fields are not generically aligned. In such a configuration the numerator of \cref{eq:OAB_energy} is expected to be strongly suppressed relative to the self-energies, so that
            \begin{equation}
                |\mathcal{O}_{AB}|\ll 1
            \end{equation}
            is the natural bound even though the topology is nontrivial.

            A simple Biot--Savart scaling makes this suppression more explicit. For thin carriers with characteristic induced speed scale \(u_0\), overlap volume \(V_{\mathrm{ov}}\), and local tangent directions \(\hat{\mathbf{t}}_A,\hat{\mathbf{t}}_B\), the cross-energy numerator scales as
            \begin{equation}
                E_{AB}^{\mathrm{int}}
                =
                \rho_{\!f}\int \mathbf{u}_A\cdot\mathbf{u}_B\,d^3x
                \;\sim\;
                \rho_{\!f}\,u_0^2\,V_{\mathrm{ov}}\,
                \big\langle \hat{\mathbf{t}}_A\cdot\hat{\mathbf{t}}_B \big\rangle_{\mathrm{ov}},
                \label{eq:hopf_scaling}
            \end{equation}
            whereas the self-energies scale as
            \begin{equation}
                E_A \sim \rho_{\!f}u_0^2V_A,
                \qquad
                E_B \sim \rho_{\!f}u_0^2V_B.
            \end{equation}
            Hence the normalized overlap obeys the reduced estimate
            \begin{equation}
                \mathcal{O}_{AB}
                \;\sim\;
                \frac{V_{\mathrm{ov}}}{\sqrt{V_A V_B}}
                \big\langle \hat{\mathbf{t}}_A\cdot\hat{\mathbf{t}}_B \big\rangle_{\mathrm{ov}}.
                \label{eq:hopf_overlap_estimate}
            \end{equation}
            For a symmetric Hopf-like or orthogonal linked configuration, \(\hat{\mathbf{t}}_A\cdot\hat{\mathbf{t}}_B\) changes sign across the overlap region and averages toward zero, so \cref{eq:hopf_overlap_estimate} makes the suppression of \(|\mathcal{O}_{AB}|\) explicit without requiring a full numerical Biot--Savart evaluation.

            Likewise, for orthogonal intersecting or weakly overlapping loops with little tangent alignment, the cross-energy integral is expected to remain small and the reduced coupling to satisfy
            \begin{equation}
                \mathcal{O}_{AB}\approx 0.
            \end{equation}
            These cases show that \(\mathcal{O}_{AB}\) is sensitive not merely to linkage or contact, but to directional and kinematic compatibility of the induced velocity fields. The present article therefore treats the coaxial examples \(\mathcal{O}_{AB}=\pm 1\) and the orthogonal/linked examples \(|\mathcal{O}_{AB}|\ll 1\) as the two simplest limiting regimes of the overlap functional.
    
    \section{Why the coherence length needs a physical decay channel}
        
        Here and below, ``entanglement-like'' continues to mean only the reduced operational notion introduced in the abstract and Introduction: nontrivial angle-dependent correlations between separated readout sites on a shared carrier.

        Suppose the two-point correlation takes the phenomenological form
        \begin{equation}
            E(\theta_A,\theta_B)
            \sim
            e^{-L/\ell_\tau}\cos(\theta_A-\theta_B).
            \label{eq:E_decay}
        \end{equation}
        Then \(\ell_\tau\) is a transport coherence length. Such a quantity requires a finite phase-memory time \(T_\phi\):
        \begin{equation}
            \ell_\tau=v_{\mathrm{tr}}\,T_\phi.
            \label{eq:l_tau_Tphi}
        \end{equation}
        At first pass, \(v_{\mathrm{tr}}=\vnorm\), so
        \begin{equation}
            \ell_\tau=\vnorm T_\phi.
        \end{equation}
        Using \(\omega_c=\vnorm/\rc\), define the phase-quality factor
        \begin{equation}
            \mathcal{Q}:=\omega_c T_\phi.
            \label{eq:Qdef}
        \end{equation}
        Then
        \begin{equation}
            \boxed{
                \ell_\tau=\rc\,\mathcal{Q}.
            }
            \label{eq:l_tau_Q}
        \end{equation}
        
        Dimensional check:
        \[
            [\ell_\tau]=[\rc][\mathcal{Q}]=\mathrm{m}.
        \]
        
        \subsection*{Interpretation}
            This is the most compact structural result of the paper: coherence length is core scale times quality factor. Once \(\omega_c\) is fixed canonically, the entire long-range correlation problem becomes a question about how large \(\mathcal{Q}\) can be.
    
    \section{Possible SST mechanisms for phase-memory loss}
        
        A finite \(\ell_\tau\) requires a specific decoherence or decorrelation mechanism. In the broader SST setting several channels are conceivable, but for definiteness the present article adopts \emph{topological leakage} as the worked reduced mechanism. The remaining channels are noted only briefly as alternatives for future analysis.
        
        \subsection{Density-noise dephasing}
            If
            \[
                \rhoF\to \rhoF+\delta\rho,
            \]
            then local phase frequencies fluctuate. A first-order estimate is
            \begin{equation}
                \frac{\delta\omega}{\omega_c}\sim \frac{\delta\rho}{\rhoF},
            \end{equation}
            so that
            \begin{equation}
                \mathcal{Q}\sim \frac{\rhoF}{\delta\rho},
                \qquad
                \ell_\tau\sim \rc\,\frac{\rhoF}{\delta\rho}.
                \label{eq:Q_density}
            \end{equation}
            This possibility is not pursued further in the present paper.
        
        \subsection{Velocity-noise dephasing}
            If
            \[
                \vnorm\to \vnorm+\delta v,
            \]
            then arrival times fluctuate from cycle to cycle. Then
            \begin{equation}
                \mathcal{Q}\sim \frac{\vnorm}{\delta v},
                \qquad
                \ell_\tau\sim \rc\,\frac{\vnorm}{\delta v}.
                \label{eq:Q_velocity}
            \end{equation}
            This possibility is likewise not pursued further here.
        
        \subsection{Topological leakage}
            If the carrier is not topologically perfect, phase information may leak into neighboring branches or through reconnection-like processes. If \(\Gamma_{\mathrm{topo}}\) denotes the corresponding leakage rate, then
            \begin{equation}
                T_\phi\sim \Gamma_{\mathrm{topo}}^{-1},
                \qquad
                \mathcal{Q}\sim \frac{\omega_c}{\Gamma_{\mathrm{topo}}},
                \qquad
                \ell_\tau\sim \frac{\vnorm}{\Gamma_{\mathrm{topo}}}.
                \label{eq:Q_topo}
            \end{equation}
            This is the mechanism adopted below for the numerical order-of-magnitude estimate, because it is the most directly tied to the persistence of shared topological carriers.
        
        \subsection{Finite clock-field coherence}
            If the foliation/clock sector has finite correlation length \(\ell_\chi\), then
            \begin{equation}
                \ell_\tau\sim \ell_\chi.
                \label{eq:l_chi}
            \end{equation}
            This would make entanglement-range a direct property of the clock field rather than only of circulation transport.
            This alternative is left for future work.

        A broader lesson of the reduced carrier picture is that \(\mathcal{Q}\) should not be read as an abstract bookkeeping symbol only. In an SST interpretation, it is more naturally understood as a measure of how long a shared circulation carrier remains topologically and dynamically coherent while staying within a resonant or quasi-resonant operating window. This interpretation is still qualitative in the present article, but it is more restrictive than a generic noise parameter and better reflects the intended role of long-lived structured carriers.
    
    \section{First numerical implications}
        
        Using
        \[
            \rc=1.40897017\times 10^{-15}\ \mathrm{m},
        \]
        we obtain
        \[
            \ell_\tau=\rc\,\mathcal{Q}
            =
            1.40897017\times 10^{-15}\,\mathcal{Q}\ \mathrm{m}.
        \]
        Hence
        \[
            \mathcal{Q}=10^6 \Rightarrow \ell_\tau\approx 1.41\times 10^{-9}\ \mathrm{m},
        \]
        \[
            \mathcal{Q}=10^{12} \Rightarrow \ell_\tau\approx 1.41\times 10^{-3}\ \mathrm{m},
        \]
        \[
            \mathcal{Q}=10^{15} \Rightarrow \ell_\tau\approx 1.41\ \mathrm{m},
        \]
        \[
            \mathcal{Q}=10^{21} \Rightarrow \ell_\tau\approx 1.41\times 10^{6}\ \mathrm{m}.
        \]
        
        Thus macroscopic or astronomical correlation ranges demand extremely large phase-quality factors. This is not an inconsistency; it is a hard closure requirement. Any successful SST entanglement sector must explain how such large \(\mathcal{Q}\) values arise, or else predict measurable finite-range deviations from ideal quantum correlations.

    \section{Worked coherence estimate from topological leakage}

        In the reduced topological-leakage closure \cref{eq:Q_topo},
        \[
            \mathcal{Q}\sim \frac{\omega_c}{\Gamma_{\mathrm{topo}}},
            \qquad
            \ell_\tau\sim \frac{\vnorm}{\Gamma_{\mathrm{topo}}}.
        \]
        These relations may be inverted to estimate the leakage rate required for a target coherence length:
        \begin{equation}
            \Gamma_{\mathrm{topo}}
            \sim
            \frac{\vnorm}{\ell_\tau}.
            \label{eq:Gammatopo_inv}
        \end{equation}

        For meter-scale coherence, \(\ell_\tau\sim 1\ \mathrm{m}\), one obtains
        \begin{equation}
            \Gamma_{\mathrm{topo}}
            \sim
            \frac{1.09384563\times 10^6\ \mathrm{m\,s^{-1}}}{1\ \mathrm{m}}
            \approx 1.09\times 10^6\ \mathrm{s^{-1}}.
        \end{equation}
        Equivalently,
        \begin{equation}
            \mathcal{Q}\sim \frac{\omega_c}{\Gamma_{\mathrm{topo}}}
            \approx
            \frac{7.76344\times 10^{20}}{1.09\times 10^6}
            \approx 7.10\times 10^{14}.
        \end{equation}
        Thus meter-scale persistence in the present toy model would require topological leakage rates no larger than order MHz. Longer coherence scales require proportionally smaller leakage rates. This does not yet establish that such values occur in SST, but it converts the phase-quality requirement into a concrete dynamical target.

        A minimal SST-style hydrodynamic grounding of this leakage rate is to parameterize it by the medium hierarchy itself. In the weak-perturbation regime one may write
        \begin{equation}
            \Gamma_{\mathrm{topo}}
            \sim
            \omega_c\,\Xi\,
            F\!\left(\frac{\rho_{\!f}}{\rhocore}\right),
            \qquad
            F(x)\to 0 \text{ as } x\to 0,
            \label{eq:Gamma_topo_general}
        \end{equation}
        where \(\Xi\) is a dimensionless distortion or reconnection factor encoding geometry, perturbation amplitude, and carrier roughness. The essential point is that topological leakage should vanish as the effective background loading becomes negligible relative to the core density.

        As a first-pass closure, the simplest monotone choice is linear:
        \begin{equation}
            \boxed{
            \Gamma_{\mathrm{topo}}
            \sim
            \omega_c\,\Xi\,\frac{\rho_{\!f}}{\rhocore}.
            }
            \label{eq:Gamma_topo_linear}
        \end{equation}
        Then
        \begin{equation}
            \mathcal{Q}
            \sim
            \frac{\omega_c}{\Gamma_{\mathrm{topo}}}
            \sim
            \frac{1}{\Xi}\frac{\rhocore}{\rho_{\!f}}.
            \label{eq:Q_density_hierarchy}
        \end{equation}
        Using the canonical SST densities
        \[
            \rhocore=3.8934358266918687\times 10^{18}\ \mathrm{kg\,m^{-3}},
            \qquad
            \rho_{\!f}=7.0\times 10^{-7}\ \mathrm{kg\,m^{-3}},
        \]
        gives
        \begin{equation}
            \frac{\rhocore}{\rho_{\!f}}
            \approx
            5.56\times 10^{24}.
            \label{eq:density_hierarchy_num}
        \end{equation}

        Using the density ratio in \cref{eq:density_hierarchy_num}, one can evaluate the natural topological-leakage limit of the linear closure \cref{eq:Gamma_topo_linear}. For a pristine shared carrier operating in a weakly perturbed regime, the geometric distortion factor is expected to be of order unity,
        \[
            \Xi \sim \mathcal{O}(1).
        \]
        The resulting intrinsic topological leakage rate is then
        \begin{equation}
            \Gamma_{\mathrm{topo}}^{\mathrm{intrinsic}}
            \sim
            \omega_c\left(\frac{\rho_{\!f}}{\rhocore}\right)
            \approx
            (7.76344\times 10^{20}\ \mathrm{s^{-1}})
            (1.7979\times 10^{-25})
            \approx
            1.40\times 10^{-4}\ \mathrm{s^{-1}}.
            \label{eq:Gamma_topo_intrinsic}
        \end{equation}

        The corresponding natural coherence length follows from the transport closure:
        \begin{equation}
            \ell_\tau^{\mathrm{intrinsic}}
            \sim
            \frac{\vnorm}{\Gamma_{\mathrm{topo}}^{\mathrm{intrinsic}}}
            \approx
            \frac{1.09384563\times 10^6\ \mathrm{m\,s^{-1}}}{1.40\times 10^{-4}\ \mathrm{s^{-1}}}
            \approx
            7.8\times 10^9\ \mathrm{m}.
            \label{eq:l_tau_intrinsic}
        \end{equation}

        Thus the SST density hierarchy, taken at face value, naturally favors extremely long phase-memory persistence. Conversely, to constrain the coherence length to the laboratory scale,
        \[
            \ell_\tau \sim 1\ \mathrm{m},
        \]
        as estimated in \cref{eq:Gammatopo_inv}, the required leakage rate
        \[
            \Gamma_{\mathrm{topo}} \approx 1.09\times 10^6\ \mathrm{s^{-1}}
        \]
        demands a very large effective distortion factor:
        \begin{equation}
            \Xi_{1\mathrm{m}}
            =
            \frac{\Gamma_{\mathrm{topo}}}{\omega_c}\left(\frac{\rhocore}{\rho_{\!f}}\right)
            \approx
            \frac{1.09\times 10^6}{7.76344\times 10^{20}}
            \left(5.56\times 10^{24}\right)
            \approx
            7.8\times 10^9.
            \label{eq:Xi_1m}
        \end{equation}

        In this reduced interpretation, laboratory-scale coherence is therefore not the natural SST baseline but a strongly distortion-dominated regime, whereas the undistorted hierarchy points toward macroscopic or even astronomical phase-memory lengths.

        Thus the SST density hierarchy by itself points toward very large \(\mathcal{Q}\), and the practical coherence length is controlled primarily by how strongly geometric distortion enhances the leakage factor \(\Xi\). In this reduced sense, the phase-memory problem becomes a competition between the enormous core-to-background density hierarchy and the geometry-dependent leakage factor \(\Xi\).

        \begin{figure}[t]
            \centering
            \begin{tikzpicture}
            \begin{axis}[
                width=0.95\linewidth,
                height=6cm,
                xlabel={effective leakage rate $\Gamma_{\mathrm{topo}}^{\mathrm{eff}}\,[\mathrm{s}^{-1}]$},
                ylabel={correlation metric},
                title={Correlation degradation versus effective leakage},
                grid=both,
                thick,
                legend pos=north east,
                legend style={font=\scriptsize, draw=none, fill=white},
            ]
            \addplot[
                blue,
                mark=*,
                mark size=1.4pt,
                solid,
            ] table [x=gamma, y=amp, col sep=comma] {fig4_leakage_curve.csv};
            \addlegendentry{correlation amplitude}
            \addplot[
                red,
                mark=square*,
                mark size=1.4pt,
                solid,
            ] table [x=gamma, y=c00, col sep=comma] {fig4_leakage_curve.csv};
            \addlegendentry{$C(0,0)$}
            \end{axis}
            \end{tikzpicture}
            \caption{Correlation degradation under an effective leakage/decoherence surrogate. The plotted quantities are the correlation amplitude and the equal-analyzer value \(C(0,0)\) versus an effective leakage rate \(\Gamma_{\mathrm{topo}}^{\mathrm{eff}}\) in the reduced model. Although this leakage term is only a toy-model surrogate rather than a first-principles SST reconnection calculation, it makes explicit how finite phase memory degrades the reduced observable correlations.}
            \label{fig:delay_phase_leakage_sweep}
        \end{figure}

        A minimal physical plausibility condition for such small \(\Gamma_{\mathrm{topo}}\) is that the shared carrier remain in a weakly perturbed, effectively inviscid regime over the observation window. In reduced SST terms, this requires at least:
        \begin{equation}
            \frac{\delta \rho_{\!f}}{\rho_{\!f}}\ll 1,
            \qquad
            \frac{\delta \vnorm}{\vnorm}\ll 1,
        \end{equation}
        together with suppression of reconnection-inducing distortions and no cavitation-like breakdown of the carrier geometry during the correlation window. In other words, the estimate \(\Gamma_{\mathrm{topo}}\sim 10^6~\mathrm{s^{-1}}\) should be read not as a universal SST prediction, but as a bound that can only be approached when the carrier remains topologically confined and dynamically quiet compared with the core-turnover scale \(\omega_c\).

    \section{Linear stability of the stationary branches}

        The stationary branches \(\Delta^\ast=n\pi\) identified in \cref{eq:stationary_branch} admit a simple linear stability analysis. Let
        \begin{equation}
            \Delta(t)=\Delta^\ast+\delta(t),
            \qquad
            |\delta(t)|\ll 1.
        \end{equation}

        For the in-phase branch \(\Delta^\ast=0\), one has \(\sin(\delta)\approx \delta\), so \cref{eq:deltaeq} linearizes to
        \begin{equation}
            \dot{\delta}(t)=-2\kappa_{AB}\,\delta(t-\tau).
            \label{eq:lin_zero}
        \end{equation}
        Seeking exponential modes \(\delta(t)\propto e^{\lambda t}\) yields the characteristic equation
        \begin{equation}
            \lambda=-2\kappa_{AB}e^{-\lambda\tau}.
            \label{eq:char_zero}
        \end{equation}
        For the scalar delay equation \(\dot{\delta}(t)=-a\,\delta(t-\tau)\) with \(a>0\), the in-phase branch is stable in the weak-delay regime when
        \begin{equation}
            a\tau<\frac{\pi}{2}.
        \end{equation}
        Hence the corresponding condition here is
        \begin{equation}
            2\kappa_{AB}\tau<\frac{\pi}{2},
            \qquad\text{equivalently}\qquad
            \mathcal{O}_{AB}\frac{L}{\rc}<\frac{\pi}{4}.
            \label{eq:weak_delay_stability}
        \end{equation}

        For the antiphase branch \(\Delta^\ast=\pi\), one has \(\sin(\pi+\delta)\approx -\delta\), so the linearized equation becomes
        \begin{equation}
            \dot{\delta}(t)=+2\kappa_{AB}\,\delta(t-\tau),
            \label{eq:lin_pi}
        \end{equation}
        with characteristic equation
        \begin{equation}
            \lambda=+2\kappa_{AB}e^{-\lambda\tau}.
            \label{eq:char_pi}
        \end{equation}

        The controlling dimensionless combination is therefore
        \begin{equation}
            2\kappa_{AB}\tau
            =
            2\,\omega_c\mathcal{O}_{AB}\frac{L}{\vnorm}
            =
            2\,\mathcal{O}_{AB}\frac{L}{\rc}.
            \label{eq:control_param}
        \end{equation}

        \begin{figure}[t]
            \centering
            \begin{tikzpicture}
            \begin{axis}[
                width=0.95\linewidth,
                height=7.0cm,
                grid=both,
                legend pos=north east,
                legend style={
                    font=\scriptsize,
                    draw=none,
                    fill=white,
                },
                xlabel={$t\,[\mathrm{s}]$},
                ylabel={$\Delta(t)\,[\mathrm{rad}]$},
                title={Representative time series across the locking sweep},
                xmin=0, xmax=8,
            ]
            \addplot[no marks, black, solid, line width=1.2pt, each nth point={80}]
                table [x=t, y=delta, col sep=comma] {fig3_locking_series_1.csv};
            \addlegendentry{$\kappa_{AB}\tau=0.35$}
            \addplot[no marks, black!55, dashdotted, line width=1.1pt, each nth point={80}]
                table [x=t, y=delta, col sep=comma] {fig3_locking_series_3.csv};
            \addlegendentry{$\kappa_{AB}\tau=2.80$}
            \addplot[no marks, black!65, loosely dashed, line width=1.0pt, each nth point={80}]
                table [x=t, y=delta, col sep=comma] {fig3_locking_series_5.csv};
            \addlegendentry{$\kappa_{AB}\tau=8.40$}
            \end{axis}
            \end{tikzpicture}
            \caption{Representative time series of the wrapped phase difference \(\Delta(t)\) across the locking sweep, shown here for three representative values of the control parameter \(\kappa_{AB}\tau\). The full-width layout emphasizes the early-time regime where the delayed trajectories differ most strongly before relaxing toward the in-phase branch.}
            \label{fig:delay_phase_locking_timeseries}
        \end{figure}

        \begin{figure}[t]
            \centering
            \begin{tikzpicture}
            \begin{axis}[
                width=0.62\linewidth,
                height=6.0cm,
                grid=both,
                legend pos=north east,
                legend style={
                    font=\scriptsize,
                    draw=none,
                    fill=white,
                },
                xlabel={$\kappa_{AB}\tau$},
                ylabel={steady-state metric},
                title={Locking diagnostics},
            ]
            \addplot[
                black,
                mark=*,
                mark size=1.2pt,
                solid,
                line width=1.0pt,
            ] table [x=ktau, y=std_delta, col sep=comma] {fig3_locking_metrics.csv};
            \addlegendentry{$\mathrm{std}(\Delta)_{\mathrm{ss}}$}
            \addplot[
                black!55,
                mark=square*,
                mark size=1.2pt,
                dashed,
                line width=1.0pt,
            ] table [x=ktau, y=dist0, col sep=comma] {fig3_locking_metrics.csv};
            \addlegendentry{$\langle |\Delta| \rangle_{\mathrm{ss}}$ to branch 0}
            \end{axis}
            \end{tikzpicture}
            \caption{Steady-state locking diagnostics versus the control parameter \(\kappa_{AB}\tau\), shown through the post-transient standard deviation \(\mathrm{std}(\Delta)_{\mathrm{ss}}\) and the mean branch distance to the in-phase branch, \(\langle |\Delta| \rangle_{\mathrm{ss}}\). This compact summary complements \cref{fig:delay_phase_locking_timeseries} by showing how the delayed branch locking changes across the sampled sweep.}
            \label{fig:delay_phase_locking_metrics}
        \end{figure}

        In the present toy-model interpretation, \cref{eq:control_param} shows that branch stability is controlled jointly by carrier geometry and overlap. The in-phase branch \(\Delta^\ast=0\) is the natural attractive branch in the weak-delay regime, while the antiphase branch \(\Delta^\ast=\pi\) is correspondingly disfavored there. A complete stability diagram for the delayed system is left for future work, but \cref{eq:char_zero,eq:char_pi,eq:control_param} identify the relevant analytic control parameter explicitly. The numerical sweep in \cref{fig:delay_phase_locking_timeseries,fig:delay_phase_locking_metrics} is included as a compact visualization of that dependence.

        It is worth noting that several macroscopic rows of Table~\ref{tab:toyparams} lie far beyond the weak-delay threshold \(\mathcal{O}_{AB}L/\rc<\pi/4\). Those rows should therefore be read only as parameter-scale illustrations, not as examples guaranteed to lie in the simple weak-delay locking regime.

    \section{Representative parameter choices and numerical regimes}

        To make the reduced model more concrete, it is useful to express the key control parameters in dimensionless form. The delayed phase system depends on
        \[
            \tau\omega_c=\frac{L}{\rc},
            \qquad
            \frac{\kappa_{AB}}{\omega_c}=\mathcal{O}_{AB},
            \qquad
            \frac{\ell_\tau}{\rc}=\mathcal{Q}.
        \]
        Hence the model is governed by three reduced quantities:
        \begin{equation}
            \Lambda:=\frac{L}{\rc},\qquad
            \mathcal{O}_{AB},\qquad
            \mathcal{Q}.
        \end{equation}

        Table~\ref{tab:toyparams} lists representative scales for several illustrative carrier lengths and overlap values. These are not claimed as unique SST predictions; they are reference regimes for reduced-model exploration.

        \begin{table}[h]
            \centering
            \caption{Representative parameter choices for the reduced delayed-phase toy model. The transport delay is \(\tau=L/\vnorm\), the coupling is \(\kappa_{AB}=\omega_c\mathcal{O}_{AB}\), and the coherence length is \(\ell_\tau=\rc\mathcal{Q}\). For illustration, a single reference value \(\mathcal{Q}=10^{15}\) is used throughout the final column.}
            \label{tab:toyparams}
            \begin{tabular}{llllll}
                \hline
                \(L\) & \(L/\rc\) & \(\tau\) & \(\mathcal{O}_{AB}\) & \(\kappa_{AB}\) & \(\ell_\tau\) for chosen \(\mathcal{Q}\) \\
                \hline
                \(10^{-12}\,\mathrm{m}\) & \(7.10\times 10^2\) & \(9.14\times 10^{-19}\,\mathrm{s}\) & \(1\) & \(7.76\times 10^{20}\,\mathrm{s^{-1}}\) & \(1.41\,\mathrm{m}\) \\
                \(10^{-9}\,\mathrm{m}\) & \(7.10\times 10^5\) & \(9.14\times 10^{-16}\,\mathrm{s}\) & \(10^{-1}\) & \(7.76\times 10^{19}\,\mathrm{s^{-1}}\) & \(1.41\,\mathrm{m}\) \\
                \(10^{-6}\,\mathrm{m}\) & \(7.10\times 10^8\) & \(9.14\times 10^{-13}\,\mathrm{s}\) & \(10^{-3}\) & \(7.76\times 10^{17}\,\mathrm{s^{-1}}\) & \(1.41\,\mathrm{m}\) \\
                \(1\,\mathrm{m}\) & \(7.10\times 10^{14}\) & \(9.14\times 10^{-7}\,\mathrm{s}\) & \(10^{-6}\) & \(7.76\times 10^{14}\,\mathrm{s^{-1}}\) & \(1.41\,\mathrm{m}\) \\
                \hline
            \end{tabular}
        \end{table}

        The main lesson of Table~\ref{tab:toyparams} is structural rather than predictive. Carrier lengths much larger than \(\rc\) correspond to extremely large reduced delays \(\Lambda=L/\rc\), while macroscopic coherence requires correspondingly large \(\mathcal{Q}\). At the same time, the effective phase-locking rate is highly sensitive to \(\mathcal{O}_{AB}\), so weak geometric overlap can strongly suppress coupling even when \(\omega_c\) is large.

    \section{Numerical illustration of the reduced delayed phase model}

        To make the reduced toy model more concrete, we include a simple numerical integration of the delayed phase system
        \begin{equation}
            \dot{\phi}_A(t)=\omega_0+\kappa_{AB}\sin\!\big[\phi_B(t-\tau)-\phi_A(t)\big],
        \end{equation}
        \begin{equation}
            \dot{\phi}_B(t)=\omega_0+\kappa_{AB}\sin\!\big[\phi_A(t-\tau)-\phi_B(t)\big],
        \end{equation}
        together with the reduced binary observables of \cref{eq:Aobs,eq:Bobs}. The purpose of these numerics is limited: they illustrate branch selection and the emergence of a nontrivial angle-dependent reduced correlation function \(C(\theta_A,\theta_B)\) within the toy model. They are not intended as a Bell/CHSH calculation.

        A minimal numerical experiment is therefore to fix \(\omega_0\), \(\tau\), and \(\kappa_{AB}\), prescribe a phase history on the interval \(t\in[-\tau,0]\), and integrate the system forward while monitoring \(\Delta(t)\) and the induced binary observables \cref{eq:Aobs,eq:Bobs}. In the absence of a complete Bell/CHSH structure, the most direct reduced-model outputs are:
        \[
            \Delta(t),
            \qquad
            A(\theta_A,t),\qquad
            B(\theta_B,t),\qquad
            C(\theta_A,\theta_B).
        \]
        In the numerical example used here, the representative integration parameters are \(\omega_0=20~\mathrm{s^{-1}}\), \(\kappa_{AB}=8~\mathrm{s^{-1}}\), \(\tau=0.35~\mathrm{s}\), \(t_{\mathrm{final}}=40~\mathrm{s}\), and a constant history \(\phi_A=0.15\), \(\phi_B=2.10\) on the interval \([-\tau,0]\).
        These values are chosen only to illustrate the dynamical behavior of the reduced mathematical model in a computationally accessible regime. They are not intended to coincide with the SST-motivated scales listed in Table~\ref{tab:toyparams}, which are many orders of magnitude more extreme.

        \begin{figure}[t]
            \centering
            \begin{tikzpicture}
            \begin{axis}[
                width=0.95\linewidth,
                height=6cm,
                xlabel={$t\,[\mathrm{s}]$},
                ylabel={$\Delta(t)=\phi_A(t)-\phi_B(t)\,[\mathrm{rad}]$},
                title={Delayed phase toy model: phase-difference time series},
                grid=both,
                thick,
                xmin=0, xmax=10,
            ]
            \addplot[
                no marks,
                very thick,
                each nth point={80},
            ] table [x=t, y=delta, col sep=comma] {fig1_delta_timeseries.csv};
            \addplot[dashed, domain=-3:3, samples=2, forget plot] ({20},{x});
            \end{axis}
            \end{tikzpicture}
            \caption{Time series of the wrapped phase difference \(\Delta(t)=\phi_A(t)-\phi_B(t)\) for a representative delayed-phase toy-model integration. After an initial transient with damped oscillations, the system relaxes toward the in-phase branch \(\Delta^\ast=0\). The dashed vertical line marks the onset of the post-transient window used for the reduced correlation analysis.}
            \label{fig:delay_phase_time_series}
        \end{figure}

        For the phase-locked branch \(\Delta^\ast=0\) and the sign-threshold observables of \cref{eq:Aobs,eq:Bobs}, one expects analytically that \(C(\theta_A,\theta_B)\) reduces to the standard triangular wave obtained by overlapping the sign of two cosine functions with relative angle shift \(\theta_B-\theta_A\). The numerical correlation plot should therefore be read as a consistency check of the reduced observable layer rather than as an independent discovery.

        To see this explicitly, take the fully locked limit \(\phi_A(t)\approx \phi_B(t)\approx \phi(t)\), and assume that the long-time dynamics sample \(\phi\) uniformly on \([0,2\pi)\). Then
        \begin{equation}
            C(\theta_A,\theta_B)
            =
            \frac{1}{2\pi}
            \int_0^{2\pi}
            \operatorname{sign}\!\big[\cos(\phi-\theta_A)\big]
            \operatorname{sign}\!\big[\cos(\phi-\theta_B)\big]
            \,d\phi.
            \label{eq:C_integral}
        \end{equation}
        By rotational invariance this depends only on \(\Delta\theta:=\theta_B-\theta_A\). Over one period, the two sign functions agree on a fraction \(1-|\Delta\theta|/\pi\) of the interval and disagree on a fraction \(|\Delta\theta|/\pi\), for \(|\Delta\theta|\le \pi\). Therefore
        \begin{equation}
            C(\Delta\theta)
            =
            \left(1-\frac{|\Delta\theta|}{\pi}\right)
            -
            \frac{|\Delta\theta|}{\pi}
            =
            1-\frac{2|\Delta\theta|}{\pi},
            \qquad
            |\Delta\theta|\le \pi,
            \label{eq:C_triangle}
        \end{equation}
        extended periodically with period \(2\pi\). This is precisely the triangular wave seen in \cref{fig:delay_phase_correlation}.

        \begin{figure}[h]
            \centering
            \begin{tikzpicture}
            \begin{axis}[
                width=0.95\linewidth,
                height=6cm,
                xlabel={$\theta_B-\theta_A\,[\mathrm{rad}]$},
                ylabel={$C(\theta_A,\theta_B)$},
                title={Reduced binary correlation versus analyzer-angle difference},
                grid=both,
                thick,
            ]
            \addplot[
                no marks,
                very thick,
            ] table [x=dtheta, y=C, col sep=comma] {fig2_correlation_curve.csv};
            \end{axis}
            \end{tikzpicture}
            \caption{Reduced binary correlation \(C(\theta_A,\theta_B)\) as a function of analyzer-angle difference \(\theta_B-\theta_A\), computed from the thresholded observables \cref{eq:Aobs,eq:Bobs} after discarding the transient. In the present representative run, the toy model yields the triangular angle dependence expected analytically for a phase-locked \(\Delta^\ast=0\) branch together with sign-threshold observables. The numerics therefore serve as a consistency check of the reduced readout layer.}
            \label{fig:delay_phase_correlation}
        \end{figure}

        Even at this reduced level, the numerical illustrations serve two immediate purposes. First, \cref{fig:delay_phase_time_series} verifies that the branch structure identified analytically is dynamically realized under explicit histories. Second, \cref{fig:delay_phase_correlation} shows that the operational observable layer \cref{eq:Aobs,eq:Bobs,eq:Cobs} produces the expected nontrivial angle dependence for representative toy-model parameters.

    
    \section{Minimal effective summary}
        
        The closure program can now be stated compactly:
        \begin{equation}
            \boxed{
                \tau=\frac{L}{\vnorm}
            }
            \label{eq:summary_tau}
        \end{equation}
        \begin{equation}
            \boxed{
                \kappa_{AB}=\omega_c\,\mathcal{O}_{AB}
            }
            \label{eq:summary_kappa}
        \end{equation}
        \begin{equation}
            \boxed{
                \ell_\tau=\rc\,\mathcal{Q}
            }
            \label{eq:summary_l}
        \end{equation}
        with
        \begin{equation}
            \omega_c=\frac{\vnorm}{\rc}.
            \label{eq:summary_omega}
        \end{equation}

        Within the toy-model hierarchy developed above, \(\mathcal{O}_{AB}\) is most naturally taken to be the energy-normalized overlap functional of \cref{eq:OAB_energy}, while \(\mathcal{Q}\) remains a reduced parameter encoding the effective phase-memory quality of the carrier. In the thin equal-length filament limit, \(\mathcal{O}_{AB}\) reduces to the simpler shared-length expression
        \[
            \mathcal{O}_{AB}=\frac{\ell_{\mathrm{shared}}}{L}.
        \]
        More generally, the intended SST reading is that \(\mathcal{O}_{AB}\) measures carrier compatibility rather than bare contact, and that \(\mathcal{Q}\) measures persistence of that compatible carrier under structured delayed transport. The long-term aim is therefore not merely to assign three free scales, but to derive them from the geometry, orientation, and dynamical quality of shared circulation structures.
    
    \section{Falsifiable consequences}
        
        The closure has three direct empirical consequences.
        Here again, ``entanglement-like'' means only the reduced operational notion introduced earlier, not a completed Bell/CHSH framework.

        The numerical illustrations are distributed accordingly across the paper: the baseline delayed locking and reduced observable behavior are shown in \cref{fig:delay_phase_time_series,fig:delay_phase_correlation}, the dependence on the control parameter \(\kappa_{AB}\tau\) is shown in \cref{fig:delay_phase_locking_timeseries,fig:delay_phase_locking_metrics}, and the effect of finite phase memory is shown in \cref{fig:delay_phase_leakage_sweep}. This placement is intended to keep each numerical figure close to the analytic discussion it supports.

        \subsection{Nontrivial angle-dependent reduced correlations}
            The binary observables \cref{eq:Aobs,eq:Bobs} define a reduced correlation function \(C(\theta_A,\theta_B)\) through \cref{eq:Cobs}. Even without a Bell analysis, the toy model predicts that phase locking on a shared carrier produces nontrivial analyzer-angle dependence whenever \(\kappa_{AB}\neq 0\). A null result \(C(\theta_A,\theta_B)\equiv 0\) across all angles would falsify the reduced readout layer of the model.
        
        \subsection{Finite-range deviation from ideal Bell scaling}
            If \(\mathcal{Q}\) is finite, then \(\ell_\tau\) is finite and ideal long-range correlations must weaken at sufficiently large \(L\):
            \begin{equation}
                E(\theta_A,\theta_B)
                \sim
                e^{-L/\ell_\tau}\cos(\theta_A-\theta_B).
            \end{equation}
            Within the present reduced-model setting, a null result out to arbitrarily large distances would force \(\mathcal{Q}\to\infty\) or invalidate the simple closure. This statement is not yet a Bell/CHSH prediction, but a modeling consequence of the finite-coherence ansatz.
        
        \subsection{Medium-sensitive dephasing}
            If \(\mathcal{Q}\) is controlled by density or velocity fluctuations, then correlation visibility should depend weakly on environmental or medium-like perturbations. A strict null result would disfavor the noise-driven versions \eqref{eq:Q_density} and \eqref{eq:Q_velocity}.
        
        \subsection{Topology-sensitive coherence}
            If \(\mathcal{Q}\) is set by topological leakage, then geometrically different production channels should carry different coherence-length budgets. This predicts channel-dependent visibility loss even at fixed separation.
    
    \section{What this does and does not solve}
        
        This article solves the dimensional and structural closure problem:
        \[
            (\tau,\kappa,\ell_\tau)
            \quad\Rightarrow\quad
            \left(
                \frac{L}{\vnorm},
                \;
                \omega_c\mathcal{O}_{AB},
                \;
                \rc\mathcal{Q}
            \right).
        \]
        It does \emph{not} yet solve the full statistical problem:
        \begin{enumerate}[label=(\arabic*)]
    \item deriving exact Bell/CHSH values,
    \item proving operational no-signaling,
    \item deriving Born probabilities from the underlying SST medium.
        \end{enumerate}
        Those remain the next closure program. What the present paper does provide is a self-contained reduced model with explicit observables, an explicit geometric coupling functional, a stability analysis, and a worked order-of-magnitude condition for maintaining long coherence. In journal-classification terms, the present manuscript should therefore be read as a dynamical-systems and reduced-modeling paper rather than a completed quantum-foundations paper.
    
    \section{Conclusion}
        
        In the present toy-model setting, the problem of delay-coupled correlations in SST can be reduced to a concrete triad of scales: delay, coupling, and coherence. Once the canonical identification
        \[
            \omega_c=\frac{\vnorm}{\rc}
        \]
        is accepted, the transport delay is fixed:
        \[
            \tau=\frac{L}{\vnorm}.
        \]
        Dimensional consistency then forces the phase-locking rate to be an intrinsic turnover frequency multiplied by a dimensionless shared-geometry factor:
        \[
            \kappa_{AB}=\omega_c\,\mathcal{O}_{AB}.
        \]
        Likewise, the correlation length is forced into the form
        \[
            \ell_\tau=\rc\,\mathcal{Q}.
        \]
        
        These relations shift the conceptual burden. In this reduced SST model, angle-dependent correlation is no longer treated purely as an abstract kinematic relation; instead, correlation strength is controlled by the degree of genuinely shared circulation structure, and correlation range by the phase-memory quality of that structure. The open problem is therefore not \emph{whether} such parameters can be written down in a toy model, but what specific SST mechanism determines \(\mathcal{O}_{AB}\) and \(\mathcal{Q}\) microscopically, and whether that mechanism can recover the stronger statistical structure of quantum theory.
        
        This is also the point at which the scope of the manuscript should be read correctly. The present article does not yet supply a Bell/CHSH derivation, a Born-rule derivation, or a direct experimental comparison. Its contribution is narrower but still substantive: it provides a reduced dynamical-systems closure in which delay, geometric coupling, and finite coherence can be parameterized self-consistently within SST.

        That is the point at which the toy model becomes sharply testable as a reduced modeling framework, even though it is not yet a full quantum-foundations theory.
            
            \begin{thebibliography}{99}
                
                \bibitem{Bell1964}
                J. S. Bell, 1964,
                \textit{On the Einstein Podolsky Rosen Paradox},
                Physics Physique Fizika \textbf{1}, 195--200.
                DOI: \href{https://doi.org/10.1103/PhysicsPhysiqueFizika.1.195}{10.1103/PhysicsPhysiqueFizika.1.195}.

                \bibitem{Yeung1999}
                M. K. S. Yeung and S. H. Strogatz, 1999,
                \textit{Time Delay in the Kuramoto Model of Coupled Oscillators},
                Physical Review Letters \textbf{82}, 648--651.
                DOI: \href{https://doi.org/10.1103/PhysRevLett.82.648}{10.1103/PhysRevLett.82.648}.
                
                \bibitem{Aspect1982}
                A. Aspect, P. Grangier, and G. Roger, 1982,
                \textit{Experimental Realization of Einstein-Podolsky-Rosen-Bohm Gedankenexperiment: A New Violation of Bell's Inequalities},
                Physical Review Letters \textbf{49}, 91--94.
                DOI: \href{https://doi.org/10.1103/PhysRevLett.49.91}{10.1103/PhysRevLett.49.91}.
                
                \bibitem{Hensen2015}
                B. Hensen et al., 2015,
                \textit{Loophole-free Bell inequality violation using electron spins separated by 1.3 kilometres},
                Nature \textbf{526}, 682--686.
                DOI: \href{https://doi.org/10.1038/nature15759}{10.1038/nature15759}.
                
                \bibitem{Erneux2009}
                T. Erneux, 2009,
                \textit{Applied Delay Differential Equations},
                Springer, New York.
                
                \bibitem{ArnoldKhesin1998}
                V. I. Arnold and B. A. Khesin, 1998,
                \textit{Topological Methods in Hydrodynamics},
                Springer, New York.
                
                \bibitem{Saffman1992}
                P. G. Saffman, 1992,
                \textit{Vortex Dynamics},
                Cambridge University Press, Cambridge.
                
                \bibitem{JacobsonMattingly2001}
                T. Jacobson and D. Mattingly, 2001,
                \textit{Gravity with a dynamical preferred frame},
                Physical Review D \textbf{64}, 024028.
                DOI: \href{https://doi.org/10.1103/PhysRevD.64.024028}{10.1103/PhysRevD.64.024028}.

                \bibitem{Madelung1927}
                E. Madelung, 1927,
                \textit{Quantentheorie in hydrodynamischer Form},
                Zeitschrift für Physik \textbf{40}, 322--326.
                DOI: \href{https://doi.org/10.1007/BF01400372}{10.1007/BF01400372}.

                \bibitem{Nelson1966}
                E. Nelson, 1966,
                \textit{Derivation of the Schrödinger Equation from Newtonian Mechanics},
                Physical Review \textbf{150}, 1079--1085.
                DOI: \href{https://doi.org/10.1103/PhysRev.150.1079}{10.1103/PhysRev.150.1079}.

                \bibitem{Barcelo2011}
                C. Barceló, S. Liberati, and M. Visser, 2011,
                \textit{Analogue Gravity},
                Living Reviews in Relativity \textbf{14}, 3.
                DOI: \href{https://doi.org/10.12942/lrr-2011-3}{10.12942/lrr-2011-3}.

                \bibitem{Volovik2003}
                G. E. Volovik, 2003,
                \textit{The Universe in a Helium Droplet},
                Oxford University Press, Oxford.
                \bibitem{Bohm1952}
                D. Bohm, 1952,
                \textit{A Suggested Interpretation of the Quantum Theory in Terms of ``Hidden'' Variables I},
                Physical Review \textbf{85}, 166--179.
                DOI: \href{https://doi.org/10.1103/PhysRev.85.166}{10.1103/PhysRev.85.166}.
            
            \end{thebibliography}

\end{document}